% GENERATED FILE. Edit canonical lessons, then run npm run generate:books.
% canonical-source-hash: fnv1a64:c00c3cda1c119c54
% canonical-lessons: ZH-C434-luyoutuan, ZH-C434-jihe, ZH-C434-gouwu, ZH-C434-liunian, ZH-C434-youji

\chapter{With a Tour Group}
\label{ch:zh-with-a-tour-group}
\hlchaptermodality{434}

\begin{chapteropening}
\textbf{By the end of this chapter:} \emph{I can talk about a tour group, where we meet, going shopping, a photo to remember and travel writing.}

Complete the last lesson of chapter 434: I can talk about a tour group, where we meet, going shopping, a photo to remember and travel writing.
\end{chapteropening}

\section[lǚyóutuán]{lǚyóutuán — tour group}

\label{lesson:ZH-C434-luyoutuan}



\begin{quote}
Before the new one: say the Mandarin for a Chinese person settled abroad, then the Mandarin for to go out to sea.
\end{quote}

\subsection*{You'll want to know: lǚyóutuán}
\emph{lǚyóutuán} — \textquotedblleft{}tour group\textquotedblright{}.

\begin{quote}
\emph{wǒ gēn yí ge lǚyóutuán zuò fēijī qù Běijīng} — I'm flying to Beijing with a tour group.
\end{quote}

Its characters come later; for now it is a word for the ear.

\begin{cousinweb}
\emph{lǚyóutuán} is built from \zh{旅游} 'to travel' and \emph{tuán} 'a group', all inherited from Middle Chinese, \emph{ljoX}, \emph{yuw} and \emph{dwan} (Guangyun, in Baxter's transcription), which glosses the last 'round'; the sense 'a group' is not given there.
\end{cousinweb}

\begin{grammarlens}[title={What you've built}]
One more word for getting around.
\end{grammarlens}

\subsection*{Your turn}
\begin{itemize}
\raggedright
  \item \emph{Say it:} \emph{lǚyóutuán}
  \item \emph{Say it:} \emph{lǚyóutuán}, once more
  \item \emph{Say it:} \emph{wǒ gēn yí ge lǚyóutuán zuò fēijī qù Běijīng} — I'm flying to Beijing with a tour group.
  \item \emph{Recall:} say the Mandarin for a Chinese person settled abroad, then the Mandarin for to go out to sea, then say \emph{lǚyóutuán} again
\end{itemize}

\begin{tcolorbox}[breakable,colback=teal!4,colframe=teal!35!black,title={Before you move on}]
What does lǚyóutuán mean? (\textquotedblleft{}Tour group\textquotedblright{}.) Say it once more.
\end{tcolorbox}
\section[jíhé]{jíhé — to gather; to assemble; to meet up as a group}

\label{lesson:ZH-C434-jihe}



\begin{quote}
Before the new one: say the Mandarin for to go out to sea, then the Mandarin for a tour group.
\end{quote}

\subsection*{You'll want to know: jíhé}
\emph{jíhé} — \textquotedblleft{}to gather; to assemble; to meet up as a group\textquotedblright{}.

\begin{quote}
\emph{míngtiān zǎoshang bā diǎn wǒmen zài jiǔdiàn dàtáng jíhé} — Tomorrow at eight in the morning we meet in the hotel lobby.
\end{quote}

Its characters come later; for now it is a word for the ear.

\begin{cousinweb}
\emph{jíhé} is built from \emph{jí} 'to gather' and \emph{hé} 'to join', both inherited from Middle Chinese, \emph{dzip} and \emph{hop} (Guangyun, in Baxter's transcription), which glosses the first 'to gather' and the second 'to join'.
\end{cousinweb}

\begin{grammarlens}[title={What you've built}]
One more word for getting around.
\end{grammarlens}

\subsection*{Your turn}
\begin{itemize}
\raggedright
  \item \emph{Say it:} \emph{jíhé}
  \item \emph{Say it:} \emph{jíhé}, once more
  \item \emph{Say it:} \emph{míngtiān zǎoshang bā diǎn wǒmen zài jiǔdiàn dàtáng jíhé} — Tomorrow at eight in the morning we meet in the hotel lobby.
  \item \emph{Recall:} say the Mandarin for to go out to sea, then the Mandarin for a tour group, then say \emph{jíhé} again
\end{itemize}

\begin{tcolorbox}[breakable,colback=teal!4,colframe=teal!35!black,title={Before you move on}]
What does jíhé mean? (\textquotedblleft{}To gather; to assemble; to meet up as a group\textquotedblright{}.) Say it once more.
\end{tcolorbox}
\section[gòuwù]{gòuwù — to go shopping; to do some shopping}

\label{lesson:ZH-C434-gouwu}



\begin{quote}
Before the new one: say the Mandarin for a tour group, then the Mandarin for to meet up as a group.
\end{quote}

\subsection*{You'll want to know: gòuwù}
\emph{gòuwù} — \textquotedblleft{}to go shopping; to do some shopping\textquotedblright{}.

\begin{quote}
\emph{xiàwǔ wǒmen zuò gōngjiāo qù shìzhōngxīn gòuwù} — In the afternoon we're taking the bus into the centre to do some shopping.
\end{quote}

Its characters come later; for now it is a word for the ear.

\begin{cousinweb}
\emph{gòuwù} is built from \emph{gòu} 'to buy' and \emph{wù} 'a thing', both inherited from Middle Chinese, \emph{kuwH} and \emph{mjut} (Guangyun, in Baxter's transcription), which glosses the first 'to buy back' and the second as the ten thousand things.
\end{cousinweb}

\begin{grammarlens}[title={What you've built}]
One more word for getting around.
\end{grammarlens}

\subsection*{Your turn}
\begin{itemize}
\raggedright
  \item \emph{Say it:} \emph{gòuwù}
  \item \emph{Say it:} \emph{gòuwù}, once more
  \item \emph{Say it:} \emph{xiàwǔ wǒmen zuò gōngjiāo qù shìzhōngxīn gòuwù} — In the afternoon we're taking the bus into the centre to do some shopping.
  \item \emph{Recall:} say the Mandarin for a tour group, then the Mandarin for to meet up as a group, then say \emph{gòuwù} again
\end{itemize}

\begin{tcolorbox}[breakable,colback=teal!4,colframe=teal!35!black,title={Before you move on}]
What does gòuwù mean? (\textquotedblleft{}To go shopping; to do some shopping\textquotedblright{}.) Say it once more.
\end{tcolorbox}
\section[liúniàn]{liúniàn — to keep as a souvenir}

\label{lesson:ZH-C434-liunian}



\begin{quote}
Before the new one: say the Mandarin for to meet up as a group, then the Mandarin for to do some shopping.
\end{quote}

\subsection*{You'll want to know: liúniàn}
\emph{liúniàn} — \textquotedblleft{}to keep as a souvenir\textquotedblright{}.

\begin{quote}
\emph{wǒmen zài ménkǒu zhàoxiàng liúniàn} — We're having our photo taken at the entrance to remember the trip.
\end{quote}

Its characters come later; for now it is a word for the ear.

\begin{cousinweb}
\emph{liúniàn} is built from \emph{liú} 'to keep' and \emph{niàn} 'to remember', both inherited from Middle Chinese, \emph{ljuw} and \emph{nemH} (Guangyun, in Baxter's transcription), which glosses the first 'to stay' and the second 'to think of'.
\end{cousinweb}

\begin{grammarlens}[title={What you've built}]
One more word for getting around.
\end{grammarlens}

\subsection*{Your turn}
\begin{itemize}
\raggedright
  \item \emph{Say it:} \emph{liúniàn}
  \item \emph{Say it:} \emph{liúniàn}, once more
  \item \emph{Say it:} \emph{wǒmen zài ménkǒu zhàoxiàng liúniàn} — We're having our photo taken at the entrance to remember the trip.
  \item \emph{Recall:} say the Mandarin for to meet up as a group, then the Mandarin for to do some shopping, then say \emph{liúniàn} again
\end{itemize}

\begin{tcolorbox}[breakable,colback=teal!4,colframe=teal!35!black,title={Before you move on}]
What does liúniàn mean? (\textquotedblleft{}To keep as a souvenir\textquotedblright{}.) Say it once more.
\end{tcolorbox}
\section[\texorpdfstring{\zh{游记} (yóujì)}{游记 (yóujì)}]{\zh{游记} (yóujì) — travel notes; account of a journey}

\label{lesson:ZH-C434-youji}



\begin{quote}
Before the new one: say the Mandarin for to do some shopping, then the Mandarin for to keep as a souvenir.
\end{quote}

\subsection*{You'll want to know: \zh{游记}}
\textbf{\zh{游记}} — \emph{yóujì} — \textquotedblleft{}travel notes; account of a journey\textquotedblright{}.

\begin{quote}
\emph{huílai yǐhòu tā xiě le yì běn yóujì} — After she came back, she wrote a book about the trip.
\end{quote}

\begin{cousinweb}
\textbf{\zh{游记}} (\emph{yóujì}) is built from \zh{游} 'to roam' and \zh{记} 'to record', both inherited from Middle Chinese, \emph{yuw} and \emph{kiH} (Guangyun, in Baxter's transcription), which glosses the second 'to record'.
\end{cousinweb}

\begin{grammarlens}[title={What you've built}]
One more word for getting around.
\end{grammarlens}

\subsection*{Your turn}
\begin{itemize}
\raggedright
  \item \emph{Say it:} \emph{yóujì}
  \item \emph{Say it:} \emph{yóujì}, once more
  \item \emph{Say it:} \emph{huílai yǐhòu tā xiě le yì běn yóujì} — After she came back, she wrote a book about the trip.
  \item \emph{Recall:} say the Mandarin for to do some shopping, then the Mandarin for to keep as a souvenir, then say \emph{yóujì} again
\end{itemize}

\begin{tcolorbox}[breakable,colback=teal!4,colframe=teal!35!black,title={Before you move on}]
What does \zh{游记} mean? (\textquotedblleft{}Travel notes; account of a journey\textquotedblright{}.) Say it once more.
\end{tcolorbox}
